\chapter{Mar.~11 --- KZ Equations}

\section{Classification of Intertwining Operators}

\begin{remark}
  Let $L_{\lambda_1}, L_{\lambda_0}$ be irreps
  of $\g$ and $k$ a complex number generic
  with respect to $\lambda_1, \lambda_0$.
  Then the induced modules
  $V_{\lambda_i, k} = \Ind_{\widetilde{g}^+}^{\widetilde{g}} L_{\lambda_i}$
  are irreducible, and
  $V_{\lambda_i, k} = L_{\lambda_i, k}$.
  Note that $V_{\lambda_i, k}[0] = L_{\lambda_i}$.

  We will be interested in $\widehat{\g}$-intertwining
  operators
  $\Phi : V_{\lambda_1, k} \to V_{\lambda_0, k}\, {\widehat{\otimes}}\, V(z)$,
  where $V(z)$ is an evaluation module and
  $\widehat{\otimes}$ denotes completed tensor product,
  i.e. elements of $V_{\lambda_0, k} \, \widehat{\otimes}\, V(z)$
  take the form
  $\sum_{i = 1}^\infty w_i \otimes v_i$, where
  $w_i \in V_{\lambda_0, k}$ are homogeneous
  with $\deg(w_i) \to -\infty$ and $v_i \in V(z)$.
  Note that we must have
  \[
    \Phi x[n] = ((x[n] \otimes 1) + z^n(1 \otimes x)) \Phi
  \]
  by the assumption that $\Phi$ is an intertwining
  operator.
\end{remark}

\begin{theorem}
  Let $g : L_{\lambda_1} \to L_{\lambda_0} \otimes V$
  be a $\g$-homomorphism. For generic $k$, there exists
  a unique $\widehat{\g}$-intertwining operator
  \[\Phi^g(z) : V_{\lambda_1, k} \longrightarrow V_{\lambda_0, k} \, \widehat{\otimes} \, V(z)\]
  such that for all $w \in V_{\lambda_1, k}[0]$,
  the degree $0$ component of $\Phi^g(z) w$ is
  equal to $gw$.
\end{theorem}

\begin{proof}
  For simplicity assume that $V$ is finite-dimensional,
  the proof for infinite-dimensional $V$ is similar.
  We construct an operator $\Phi^g(z)$ satisfying the
  given conditions uniquely.
  Let $w \in V_{\lambda_1, k}[0] = L_{\lambda_1}$.
  Then $\Phi^g(z) w$ is annihilated by
  $\g \otimes t \C[t]$ and the degree $0$ component
  of $\Phi^g(z) w$ is $gw$. Thus
  \[
    \Phi^g(z) w
    \in (V_{\lambda_0, k} \, \widehat{\otimes} \, V(z))^{\g \otimes t \C[t]}
    \cong \Hom_{\g \otimes t \C[t]}(V_{\lambda_0, k}^*, V(z)).
  \]
  Note that $V_{\lambda_0, k}^*$ is freely
  generated by $(V_{\lambda, k}[0])^* = L_{\lambda_0}^*$
  over $\g \otimes t \C[t]$ for generic $k$, so
  $\Phi^g(z) w$ is uniquely determined by its
  restriction in
  $\Hom_\C(L_{\lambda_0}^*, V) \cong L_{\lambda_0} \otimes V$, i.e.
  its degree $0$ component.
  Using
  \[
    \Phi^g(z) x[n] w = ((x[n] \otimes 1) + z^n(1 \otimes x)) \Phi^g(z) w
  \]
  for $w \in V_{\lambda_1, k}$ and $n < 0$,
  we can extend $\Phi^g(z)$ to all of $V_{\lambda_1, k}$
  inductively, as $V_{\lambda_1, k}$ is generated by
  its degree $0$ component over
  $\g \otimes t^{-1} \C[t^{-1}]$.
  This defines the operator $\Phi^g(z)$.
  Now note that no choices were allowed in the
  above construction, hence $\Phi^g(z)$ must be unique.
\end{proof}

\begin{prop}
  For any $\Delta \in \C$, the expression
  $z^{-\Delta} \Phi^g(z)$ is also a
  $\widehat{\g}$-intertwiner
  \[V_{\lambda_1, k} \longrightarrow V_{\lambda_0, k} \, \widehat{\otimes} \, z^{-\Delta} V[z, z^{-1}].\]
  It is a $\widetilde{\g}$-homomorphism if and only if
  $\Delta = \Delta(\lambda_1) - \Delta(\lambda_0)$.
\end{prop}

\begin{proof}
  Note that $z^{-\Delta} \Phi^g(z)$ is a
  $\widetilde{\g}$-homomorphism if and only if
  \[
    z^{-\Delta} \Phi^g(z) d
    - (d \otimes 1) z^{-\Delta} \Phi^g(z)
    = \bigg(1 \otimes z \frac{d}{dz}\bigg)
    z^{-\Delta} \Phi^g(z).
  \]
  Since $z^{-\Delta} \Phi^g(z) = \sum_{n \in \Z} \Phi^g[n] z^{-n - \Delta}$
  and $\Phi^g[n]$ has degree $n$, it suffices
  to have the above identity in degree $0$,
  which is given by $-\Delta(\lambda_1) + \Delta(\lambda_0) = - \Delta$.
\end{proof}

\begin{remark}
  Denote $\widetilde{\Phi}^g(z) = z^{-\Delta} \Phi^g(z) = \sum_{n \in \Z} \Phi^g[n] z^{-n - \Delta}$
  for $\Delta = \Delta(\lambda_1) - \Delta(\lambda_0)$.
\end{remark}

\section{Operator KZ Equations}

\begin{remark}
  Let $V^*$ be the restricted dual of $V$, so that
  $V^*(z) \cong (V(z))^*$ as $\widetilde{\g}$-modules.
  Let $u \in V^*$, and
  define $\Phi^g_u(z) w = u(\Phi^g(z) w)$
  for $w \in V_{\lambda_1, k}$,
  so that $\Phi^g_u(z)$ is an operator
  $V_{\lambda_1, k} \to \widehat{V}_{\lambda_0, k}$.
  The intertwining relation is given by
  $[x[n], \Phi^g_u(z)] = z^n \Phi^g_{xu}(z)$.
  Note that the currents act on $V^*(z)$ by
  \begin{gather*}
    J_x^+(\zeta) u
    = \sum_{n > 0} \zeta^{n - 1} x[-n] u
    = \sum_{n > 0} \zeta^{n - 1} z^{-n} xu
    = \frac{xu}{z - \zeta}, \quad |\zeta| < |z|, \\
    J_x^-(\zeta) u
    = - \sum_{n \le 0} \zeta^{n - 1} x[-n] u
    = - \sum_{n \le 0} \zeta^{n - 1} z^{-n} xu
    = \frac{xu}{z - \zeta}, \quad |z| < |\zeta|,
  \end{gather*}
  and in terms of currents, the intertwining
  relation is given by
  \[
    [J_x^{\pm}(\zeta), \Phi^g_{u}(z)]
    = \frac{1}{z - \zeta} \Phi^g_{xu}(z). \tag{$*$}
  \]
  Let $V = V_\mu$ be a lowest weight $\g$-module
  with lowest $-\mu$, so that the Casimir element
  $C$ acts on $V_\mu$ by the scalar $(\mu, \mu + 2\rho)$.
  Let $\widehat{\Phi}^g(z) = z^{-\Delta(\mu)} \widetilde{\Phi}^g(z) = \sum_{n \in \Z} \Phi^g[n] z^{-n - \Delta}$
  for $\Delta = \Delta(\lambda_1) - \Delta(\lambda_0) + \Delta(\mu)$.


  For $a \in \g$, define the \emph{normal ordered product}
  \[
    {:} J_a(z) \widehat{\Phi}^g_u(z) {:}
    = J_a^+(z) \widehat{\Phi}^g_u(z)
    - \widehat{\Phi}^g_u(z) J_a^-(z).
  \]
  The following result is known as the
  \emph{operator Knizhnik-Zamolodchikov (KZ) equation}.
\end{remark}

\begin{theorem}[Frenkel-Reshetikhin]\label{thm:operator-KZ}
  Let $B$ be an orthonormal basis for $\g$. Then
  \[
    (k + h^\vee) \frac{d}{dz} \widehat{\Phi}^g_u(z)
    = \sum_{a \in B} {:} J_a(z) \widehat{\Phi}^g_{au}(z) {:}.
  \]
\end{theorem}

\begin{proof}
  Since $\widetilde{\Phi}^g(z)$ is invariant
  under $d$, we have
  $z (d / dz) \widetilde{\Phi}^g_u(z) = -[d, \widetilde{\Phi}^g_u(z)]$,
  or equivalently,
  \[
    z \frac{d}{dz} \widehat{\Phi}^g_u(z)
    = -[d, \widehat{\Phi}^g_u(z)]
    - \Delta(\mu) \widehat{\Phi}^g_u(z).
  \]
  Using the expansion for $d = -L_0$, the above
  equality becomes
  \[
    z \frac{d}{dz} \widehat{\Phi}^g_u(z)
    = \frac{1}{2(k + h^\vee)} \sum_{a \in B}
    \bigg(\underbrace{\sum_{n \le 0} [a[n] a[-n], \widehat{\Phi}^g_u(z)] + \sum_{n > 0} [a[-n] a[n], \widehat{\Phi}^g_u(z)]}_{S(a)}\bigg)
    - \Delta(\mu) \widehat{\Phi}^g_u(z).
  \]
  By the intertwining relation $(*)$, we can
  rewrite the inner sum above as
  \begin{align*}
    S(a)
    &= \sum_{n \le 0} (z^{-n} a[n] \widehat{\Phi}^g_{au}(z) + z^n \widehat{\Phi}^g_{au}(z) a[-n])
    + \sum_{n > 0} (z^n a[-n] \widehat{\Phi}^g_{au}(z) + z^{-n} \widehat{\Phi}^g_{au}(z) a[n]) \\
    &= 2 z J_a^+(z) \widehat{\Phi}^g_{au}(z)
    - 2z \widehat{\Phi}^g_{au}(z) J_a^-(z)
    + a[0] \widehat{\Phi}^g_{au}(z) - \widehat{\Phi}^g_{au}(z) a[0]
    = 2 z {:} J_a(z) \widehat{\Phi}^g_{au}(z) {:}
    + \widehat{\Phi}^g_{a^2 u} (z).
  \end{align*}
  Substituting this back in, we get
  \[
    z \frac{d}{dz} \widehat{\Phi}^g_u(z)
    = \frac{1}{k + h^\vee} \sum_{a \in B}
    \bigg(z {:} J_a(z) \widehat{\Phi}^g_{au}(z) {:} + \frac{1}{2} \widehat{\Phi}^g_{a^2 u}(z)\bigg)
    - \Delta(\mu) \widehat{\Phi}^g_u(z).
  \]
  Now $\sum_{a \in B} \widehat{\Phi}^g_{a^2 u}(z) = \widehat{\Phi}^g_{C u}(z) = (\mu, \mu + 2 \rho) \widehat{\Phi}^g_u(z)$,
  so the last two terms cancel, and we have
  \[
    z \frac{d}{dz} \widehat{\Phi}^g_u(z)
    = \frac{1}{k + h^\vee} \sum_{a \in B} z {:} J_a(z) \widehat{\Phi}^g_{au}(z) {:}.
  \]
  Rearranging gives the desired equation from here.
\end{proof}

\begin{prop}
  We have the following:
  \[
    [L_m, \widehat{\Phi}^g_u(z)]
    = \frac{z^m}{k + h^\vee} {:} J(z) \widehat{\Phi}^g_u(z) {:}
    + m z^m \Delta(\mu) \widehat{\Phi}^g_u(z).
  \]
\end{prop}
